\subsubsection{Inter-Class Variability}\label{sec:inter-class-variability}

This section introduces another fundamental challenge that is \textbf{different from all the previous ones}. Until now, we always assumed: ``\emph{The object is the same, but the image changes}''. Now the situation is reversed: ``\emph{The label is the same, but the objects themselves can be very different}''.

\highspace
The fourth challenge in image classification is \definition{Inter-Class Variability}. The idea is very simple: \textbf{images belonging to the same class may look dramatically different from one another}. Previously we had translation, rotation, illumination, deformation, viewpoint, where \textbf{one object} generated many different images. Now we have \textbf{many different objects} that all belong to \textbf{the same semantic class}.

\highspace
\textcolor{Green3}{\faIcon{question-circle} \textbf{What does ``inter-class variability'' mean?}} The name can be a little confusing. Here, ``\emph{variability}'' refers to the \textbf{variation inside a semantic class}. For example, consider the class ``dog''. This class includes, Labradors, Bulldogs, Golder Retrievers, Epagneul Breton, which may have completely different colors, sizes, ears, body portions, and so on. Yet every image must receive exactly the same label: ``dog''.

\highspace
\begin{flushleft}
    \textcolor{Red2}{\faIcon{exclamation-triangle} \textbf{Why is inter-class variability a problem?}}
\end{flushleft}
Recall our linear classifier:
\begin{equation*}
    \mathbf{s} = W\mathbf{x} + b
\end{equation*}
Earlier  interpreted every row of $W$ as a \textbf{template} for a class. For example, template for class ``dog'' would be a prototypical dog image. But now, \textbf{\emph{what should the template for class ``dog'' look like?}} Should it be a Labrador, a Bulldog, or a Golden Retriever? The answer is: \textbf{there is no single appearance that represents every possible dog}. A single template therefore struggles to match all the visual variations inside the class. The humans do not memorize every dog individually. Instead, we learn the abstract concept of ``dog'' and can recognize a dog even if we have never seen that particular dog before. Deep neural networks attempt to learn these higher-level concepts automatically.

\highspace
\begin{flushleft}
    \textcolor{Red2}{\faIcon{question-circle} \textbf{Why data augmentation does not solve inter-class variability?}}
\end{flushleft}
Earlier we said that data augmentation helps because we generate new images of the same object. But can data augmentation help with inter-class variability? The answer is \textbf{no}. Data augmentation can only generate new images of the same object, but it cannot generate new objects that belong to the same class. Those are \textbf{different objects}, not different transformations of the same object. So the \hl{model still needs many real examples from the class to learn all its variations.}

\highspace
\textcolor{Green3}{\faIcon{check-circle} \textbf{Why CNNs help with inter-class variability?}} CNNs do not learn one giant template for each class (e.g., for a ``dog''). Instead, they learn many reusable features, such as ``ear'', ``tail'', ``eye'', ``nose'', and so on. These features can be combined in many different ways to recognize many different objects that belong to the same class.